\documentclass[12pt]{article}
\usepackage[utf8]{inputenc}
\usepackage[T1]{fontenc}
\usepackage{hyperref}
\title{Northern Tlaxiaco Mixtec Language Resources for K-12 Curriculum Development}
\author{Compiled by Claude Sonnet 5.5 (claude-sonnet-5-5) for Andr\'as Kornai}
\date{}
\begin{document}
\maketitle
\section{Introduction}
Northern Tlaxiaco Mixtec (ISO 639-3 xtn, Glottocode nort2985) is an Otomanguean language of Oaxaca, Mexico, with Spanish as the secondary language. The search found no speaker counts. Materials found are mostly Bible translations and audio, plus two older SIL publications cited by ASJP.
\section{Primary Educational Resources}
No primers, textbooks or school programmes were found in the search results.
\section{Linguistic and Reference Materials}
\begin{itemize}
\item ASJP entry, which cites two SIL (Instituto Ling\"u\'istico de Verano) publications on the orthography and sounds of the language (2005 and 2003): \url{https://asjp.clld.org/languages/MIXTEC_NORTHERN_TLAXIACO}
\end{itemize}
\section{Religious and Community Resources}
\begin{itemize}
\item New Testament in Northern Tlaxiaco Mixtec, Digital Bible Society: \url{https://dbs.org/bibles/XTNWBT}
\item eBible.org edition (Wycliffe, all rights reserved, so check terms before reuse): \url{https://ebible.org/xtn/copyright.htm}
\item Joshua Project profile: \url{https://www.joshuaproject.net/index.php/people_groups/14036/US}
\end{itemize}
\section{Digital and Audio Resources}
\begin{itemize}
\item Global Recordings Network, Mixtec, Northern Tlaxiaco (Bible lessons in 40 sections, MP3 and MP4 packages): \url{https://www.globalrecordings.net/en/language/xtn}
\item Global Recordings, Yoso\~nama variety: \url{https://www.globalrecordings.net/en/language/14129}
\item Global Recordings, Santiago Nundiche variety (may have no audio yet): \url{https://glorecnet-website.k.globalrecordings.net/en/language/25422}
\end{itemize}
\section{Curriculum Development Context}
A written standard exists in the New Testament and the SIL orthography publications, which gives a starting point. No school materials were found, and variety differences would need attention.
\end{document}